% viveka-insight system paper
% Compile: pdflatex viveka_insight_paper.tex (twice)
% TODO before submission: adapt to the target journal's class file (this
% draft uses plain `article` for portability). Optionally name the
% volunteer annotator in the Acknowledgements, with their consent.
\documentclass[11pt,a4paper]{article}

\usepackage[utf8]{inputenc}
\usepackage[T1]{fontenc}
\usepackage{lmodern}
\usepackage{microtype}
\usepackage[margin=2.5cm]{geometry}
\usepackage{booktabs}
\usepackage{graphicx}
\usepackage{tikz}
\usetikzlibrary{arrows.meta,positioning,shapes.geometric}
\usepackage{xcolor}
\usepackage{enumitem}
\usepackage[hidelinks]{hyperref}
\usepackage{url}

\newcommand{\system}{\textsc{Viveka-Insight}}

\title{\system{}: A Cross-Lingual Concept-Graph Retrieval System and
Citation-Grounded Question Answering over the Complete Works of
Swami Vivekananda}

\author{%
  Tamal Maharaj\thanks{Corresponding author: \texttt{tamal@gm.rkmvu.ac.in}.
  ORCID 0009-0001-5835-8967.}\\
  \small Ramakrishna Mission Vivekananda Educational and Research Institute\\
  \small Belur Math, Howrah, West Bengal, India
}
\date{July 2026}

\begin{document}
\maketitle

\begin{abstract}
Classical spiritual and philosophical corpora pose three compounding
challenges for modern retrieval systems: they exist in multiple languages
without parallel alignment, their vocabulary differs sharply from the
vocabulary of contemporary readers, and generated answers over such
culturally sensitive material must be verifiably grounded in the source
text. We present \system{}, an end-to-end open system that addresses all
three for the works of Swami Vivekananda (1863--1902): the nine-volume
English \emph{Complete Works} and the ten-volume Bengali
\emph{Vani o Rachana}, two related but non-parallel corpora totalling
about 15 million characters. The system combines (i) multi-granular dense
retrieval at sentence, paragraph, and chapter level using a single
multilingual embedding space; (ii) an LLM-extracted \emph{concept graph}
of 8{,}362 language-agnostic concepts with 87{,}499 paragraph--concept and
55{,}865 concept--concept edges, which serves as a shared abstraction
layer bridging the two languages; and (iii) a question-answering pipeline
that explicitly \emph{bridges modern questions} (e.g., ``how do I get rid
of mobile phone addiction?'') to era-appropriate concepts such as
\emph{attachment} and \emph{self-control} before retrieval, then generates
answers in which every claim carries a citation deep-linked to the exact
source paragraph, with unsupported citation markers removed. The full index is
built for under one GPU-day and serves interactive queries in 0.3 seconds
and cited answers in 6--9 seconds on a single GPU. In a known-item
cross-lingual evaluation over 194 automatically verified
English--Bengali rendered lecture pairs, the production configuration
retrieves the counterpart chapter at Recall@10 of 0.86 in both
directions; over a 30-question audit, the bridging stage fires on 93\% of
modern questions versus 13\% of timeless ones, and no generated citation
marker pointed outside the retrieved evidence. A human study of 200
randomly sampled paragraph--concept edges places concept-extraction
precision at 0.60 under the strict consensus of two annotators
(Cohen's $\kappa = 0.61$), and shows that the extractor's own confidence
weight is calibrated: restricting the graph to weight $\geq 0.8$ raises
consensus precision to 0.71 while still covering 98\% of the paragraphs
that carry any concept. Precision is markedly lower on Bengali than on
English material (0.54 vs 0.68), locating the concept layer's weakness
in exactly the half of the corpus that cross-lingual bridging most
depends on. We describe the
architecture, report these studies alongside corpus and graph statistics,
and present case studies of temporal concept bridging in both English and
Bengali. The design generalizes to other multilingual classical corpora.
\end{abstract}

\section{Introduction}

Large digitized editions of classical religious and philosophical
literature are now widely available, but access to them remains largely
lexical: readers search for words, while what they hold are questions.
The gap is widest for corpora like the works of Swami Vivekananda, a
central figure of modern Vedanta whose collected English works span nine
volumes and whose Bengali writings and recorded conversations fill ten
more. The two collections overlap thematically but are \emph{not}
translations of one another; a reader fluent in only one language has no
lexical route into roughly half of the material.

A second, subtler gap is \emph{temporal}. Contemporary readers bring
contemporary questions --- about smartphone addiction, social-media envy,
or career anxiety --- to an author who died in 1902. The vocabulary of
such questions simply does not occur in the corpus, so even strong dense
retrievers underperform; yet the corpus deals extensively with the
underlying phenomena under names such as \emph{attachment},
\emph{control of the mind}, and \emph{habit}. Bridging this gap is not a
retrieval nicety but the central requirement for making a
nineteenth-century corpus useful to twenty-first-century readers.

Finally, any system that \emph{generates} answers about revered texts
must confront hallucination \cite{ji2023survey}. Attributing invented
statements to a religious figure is a materially worse failure mode than
a wrong answer in open-domain QA, so verifiable grounding --- every claim
traceable to a specific passage a reader can open --- is a hard
constraint rather than a desirable extra.

\system{} is a complete, deployed system built around these three
requirements. Its contributions are:

\begin{enumerate}[leftmargin=1.4em]
  \item \textbf{A cross-lingual concept graph as retrieval
  infrastructure.} An instruction-tuned LLM reads every paragraph of both
  corpora and emits concepts with \emph{canonical English, lowercase,
  hyphenated labels}. Because the label space is shared, a Bengali
  paragraph about \emph{vairagya} and an English lecture on renunciation
  attach to the same node; string equality on labels gives cross-lingual
  linking without any parallel data (Section~\ref{sec:graph}).
  \item \textbf{Multi-granular fused retrieval.} Queries are matched
  against sentence-, paragraph-, and chapter-level indexes in both
  languages \emph{and} against the concept graph; candidates are fused
  with weighted reciprocal rank fusion \cite{cormack2009rrf} and reranked
  by a multilingual cross-encoder (Section~\ref{sec:retrieval}).
  \item \textbf{Temporal concept bridging for QA.} A dedicated analysis
  step maps modern questions onto the concept vocabulary that actually
  exists in the graph before retrieval, and the mapping is surfaced to
  the user as an explicit, inspectable ``bridge''
  (Section~\ref{sec:qa}).
  \item \textbf{Citation-grounded generation.} Answers are composed only
  from retrieved passages; each claim carries a bracketed citation that
  is rewritten into a deep link to the exact \emph{paragraph} of the
  published source, and citation markers that do not correspond to a
  retrieved passage are deleted (Section~\ref{sec:qa}).
  \item \textbf{An economical, reproducible build.} The full pipeline
  runs on a single GPU in under a day, incremental rebuilds after source
  edits take 10--15 minutes via a warm-start cache, and the system serves
  interactive latencies on modest hardware (Section~\ref{sec:stats}).
\end{enumerate}

\section{Related Work}

\textbf{Dense retrieval and RAG.} Dense passage retrieval
\cite{karpukhin2020dpr} and retrieval-augmented generation
\cite{lewis2020rag,gao2023ragsurvey} are the standard architecture for
grounded QA. We adopt the general recipe but depart in two ways: the
retrieval layer is multi-granular and graph-augmented, and generation is
constrained to a citation discipline enforced by post-processing rather
than trust in the model.

\textbf{Graph-augmented retrieval.} GraphRAG \cite{edge2024graphrag}
builds an entity graph over a corpus with an LLM and uses community
summaries for query-focused summarization. \system{} shares the
LLM-builds-a-graph idea but uses the graph differently: nodes are
\emph{concepts} with language-agnostic canonical labels, and the graph's
primary role is cross-lingual bridging and query-time expansion into two
non-parallel corpora, rather than corpus summarization.

\textbf{Multilingual embeddings.} We build on BGE-M3
\cite{chen2024bgem3}, whose single embedding space covers both English
and Bengali and supports inputs long enough for chapter-level
representation, and on its companion multilingual cross-encoder reranker.
Sentence-level transformer embeddings follow the general approach of
\cite{reimers2019sbert}.

\textbf{QA over religious corpora.} Verse-based test collections and QA
systems exist for the Qur'an \cite{malhas2020ayatec} and other
scriptures; these typically operate in a single language over a closed
canon. To our knowledge, \system{} is the first system to combine
non-parallel bilingual retrieval, an LLM-built concept layer, and
explicit modern-question bridging for a classical corpus.

\section{Corpus and Preprocessing}
\label{sec:corpus}

The English source is the nine-volume \emph{Complete Works of Swami
Vivekananda} (lectures, letters, recorded conversations, and writings;
$\sim$8.1 million characters). The Bengali source is the ten-volume
\emph{Vani o Rachana} ($\sim$6.9 million characters), which contains
original Bengali writings and conversations as well as independent
Bengali renderings of some English lectures; no paragraph-level alignment
between the two exists, and we do not attempt to create one.

Both corpora are parsed from their published HTML editions with
structure-aware parsers that recover the volume $\to$ (section) $\to$
chapter $\to$ paragraph hierarchy and preserve both the chapter and the
per-paragraph HTML anchor identifiers, which later enable paragraph-precise
deep-linking of citations back into the publicly hosted editions. Sentences are split with language-appropriate
rules (NLTK \cite{bird2004nltk} for English; a rule-based splitter
honouring the \emph{danda} and Bengali punctuation for Bengali). Bengali
text is NFC-normalized throughout; a single normalization function is the
sole source of truth for text keys across the pipeline, which the
incremental-rebuild machinery (Section~\ref{sec:build}) depends on.
Table~\ref{tab:corpus} summarizes the corpus.

\begin{table}[t]
\centering
\caption{Corpus statistics after parsing.}
\label{tab:corpus}
\begin{tabular}{lrrr}
\toprule
 & English & Bengali & Total\\
\midrule
Volumes (published)  & 9      & 10     & 19\\
Chapters    & 1{,}452 & 539   & 1{,}991\\
Paragraphs  & 18{,}239 & 14{,}455 & 32{,}694\\
Sentences   & 88{,}150 & 80{,}692 & 168{,}842\\
Characters  & 8.1M   & 6.9M   & 15.0M\\
\bottomrule
\end{tabular}
\end{table}

\section{The Concept Graph}
\label{sec:graph}

\subsection{LLM concept extraction}

An instruction-tuned LLM (Qwen2.5-14B-Instruct \cite{qwen25}, served with
vLLM \cite{kwon2023vllm}) reads every paragraph in both languages and
emits structured JSON containing (i) a one-line English summary, (ii) a
set of \emph{concepts}, and (iii) a set of named \emph{entities} typed as
person, place, text, or deity. The extraction prompt imposes the single
most important invariant of the system: \textbf{concept labels are
canonical English, lowercase, and hyphenated} regardless of the paragraph
language. The canonical label thus acts as a cross-lingual key: the
Bengali surface form and the English surface form of a concept are stored
as \emph{aliases} (60{,}838 aliases; 30{,}051 English, 30{,}787 Bengali),
while identity lives in the label itself. Each paragraph--concept edge
carries a relation (\emph{discusses}, \emph{exemplifies}, or
\emph{contrasts}) and a weight.

Near-duplicate concepts produced by paraphrase (e.g., spelling variants)
are merged by embedding the labels and clustering above a conservative
cosine threshold (0.92). Two kinds of concept--concept edges are then
built: \emph{similar} edges (label-embedding cosine $\geq$ 0.78, top-12
neighbours per node) and \emph{co-occurs} edges (concepts extracted from
the same paragraph at least twice corpus-wide).
Table~\ref{tab:graph} summarizes the resulting graph;
Table~\ref{tab:topconcepts} lists the most frequent concepts, which
align well with the corpus's known thematic centres.

\begin{table}[t]
\centering
\caption{Concept-graph statistics.}
\label{tab:graph}
\begin{tabular}{lr}
\toprule
Concepts (canonical labels) & 8{,}362\\
Concept aliases (EN / BN) & 30{,}053 / 30{,}797\\
Entities (person / other / place / text / deity) & 3{,}688 / 2{,}372 / 1{,}275 / 846 / 492\\
Paragraph$\to$concept edges & 87{,}518\\
\quad discusses / exemplifies / contrasts & 75{,}403 / 8{,}735 / 3{,}380\\
Concept$\leftrightarrow$concept edges & 55{,}872\\
\quad similar / co-occurs & 29{,}482 / 26{,}390\\
Paragraphs with $\geq$1 concept & 23{,}363 (71.5\%)\\
\quad retained at extractor weight $\geq$ 0.8 & 22{,}914 (98.1\%)\\
\bottomrule
\end{tabular}
\end{table}

\begin{table}[t]
\centering
\caption{Most frequent concepts by paragraph mentions.}
\label{tab:topconcepts}
\begin{tabular}{lr@{\qquad}lr}
\toprule
Concept & Mentions & Concept & Mentions\\
\midrule
self-realization      & 4{,}521 & renunciation & 2{,}553\\
knowledge-realization & 3{,}858 & liberation   & 2{,}528\\
devotion              & 3{,}179 & patience     & 1{,}787\\
maya                  & 2{,}684 & duty         & 1{,}437\\
non-attachment        & 2{,}669 & karma        & 1{,}334\\
\bottomrule
\end{tabular}
\end{table}

The 71.5\% paragraph coverage is deliberate rather than a deficiency:
short fragments (single-word replies in recorded conversations, headings,
brief exclamations) legitimately carry no concept content, and forcing
extraction on them degrades graph precision.

Coverage, however, is not correctness. A two-annotator study of the
edges themselves (Section~\ref{sec:eval-precision}) puts their
strict-consensus precision at 0.60, rising to 0.71 on the
high-confidence subgraph and falling to 0.54 on Bengali material; the
graph should be read as a broad, imperfect expansion layer rather than a
curated index, and everything we claim for it below is calibrated to
that.

\subsection{Economical construction and incremental rebuilds}
\label{sec:build}

Concept extraction is the only expensive stage ($\sim$3.5 hours for the
full corpus on a single 80\,GB-class GPU); everything else (parsing,
embedding, graph assembly) takes minutes. To make the corpus
\emph{editable} --- OCR fixes and structural cleanups are routine in
digitized classical texts --- the pipeline snapshots all extraction
outputs keyed by (NFC-normalized paragraph text, language, extractor
version) before any re-parse, and restores them onto the fresh paragraph
rows afterwards. Only genuinely new or changed paragraphs reach the LLM.
An incremental rebuild after a content edit completes in 10--15 minutes,
which in practice is the difference between a maintained resource and a
frozen one.

\section{Multi-Granular Fused Retrieval}
\label{sec:retrieval}

\begin{figure}[t]
\centering
\definecolor{figgrey}{HTML}{BFBFBF}
\definecolor{figgreen}{HTML}{D9EAD3}
\definecolor{figblue}{HTML}{CFE2F3}
\definecolor{figyellow}{HTML}{FCE7C3}
\definecolor{figpink}{HTML}{F7D2D5}
\resizebox{\textwidth}{!}{%
\begin{tikzpicture}[
  font=\small,
  box/.style={align=center, inner sep=7pt, minimum height=13mm},
  arr/.style={-{Stealth[length=3mm,width=2.6mm]}, line width=1.1pt,
              draw=gray!75, rounded corners=6pt}
]
\node[box, fill=figgrey, minimum width=26mm] (q) at (0,0)
  {Query\\(EN or BN)};
\node[box, fill=figgreen, minimum width=52mm, minimum height=17mm]
  (vec) at (5.4,2.1) {Sentence/Paragraph/Chapter\\
  FAISS indexes $\times$ \{EN, BN\}};
\node[box, fill=figgreen, minimum width=52mm, minimum height=17mm]
  (con) at (5.4,-2.1) {Concept Index +\\
  1-hop graph-walk $\rightarrow$\\ Linked Paragraphs};
\node[box, fill=figblue, minimum width=30mm] (rrf) at (10.6,0)
  {Weighted RRF\\Fusion};
\node[box, fill=figyellow, minimum width=34mm] (rer) at (14.4,0)
  {Cross Encoder\\rerank (Top 60)};
\node[box, fill=figpink, minimum width=46mm] (out) at (18.9,0)
  {Ranked Passages +\\\emph{via-concept} explanations};
\draw[arr] (q.east)   -- ++(0.45,0) |- (vec.west);
\draw[arr] (q.east)   -- ++(0.45,0) |- (con.west);
\draw[arr] (vec.east) -- ++(0.45,0) |- (rrf.west);
\draw[arr] (con.east) -- ++(0.45,0) |- (rrf.west);
\draw[arr] (rrf.east) -- (rer.west);
\draw[arr] (rer.east) -- (out.west);
\end{tikzpicture}%
}
\caption{Retrieval architecture. Path A (top) matches the query against
three granularities per language; Path B (bottom) routes it through the
concept graph. Both feed a single fusion and reranking stage.}
\label{fig:arch}
\end{figure}

All text units --- 168k sentences, 33k paragraphs, 2k chapters, and the
8.4k concept labels --- are embedded into the same 1024-dimensional
space with BGE-M3 \cite{chen2024bgem3} and stored in exact
inner-product FAISS indexes \cite{johnson2019faiss} (about 1.1\,GB in
total including metadata, small enough that approximate indexing is
unnecessary). Retrieval runs two paths per language
(Figure~\ref{fig:arch}):

\textbf{Path A (direct).} The query is matched against the sentence
(top-30), paragraph (top-20), and chapter (top-10) indexes. Sentence and
chapter hits are mapped to their enclosing/leading paragraphs, since the
paragraph is the unit surfaced to users; the granularity that found each
hit is retained and displayed.

\textbf{Path B (concept-mediated).} The query is matched against the
concept-label index (top-12 anchors), the anchors are expanded one hop
along \emph{similar} and \emph{co-occurs} edges with weight decay, and
each activated concept contributes its top linked paragraphs. Because
concept nodes are language-agnostic, an English query activates Bengali
paragraphs (and vice versa) with no translation step. Path B also yields
an explanation --- the concepts through which each paragraph was reached
--- which the interface displays as ``via concepts'' badges.

Candidates from the four rank lists (three granularities + concept path)
are fused with weighted reciprocal rank fusion \cite{cormack2009rrf},
with the concept path weighted slightly above the direct paths (1.2 vs.\
1.0) and chapters below (0.7). The top 60 fused candidates are reranked
by the BGE-reranker-v2-m3 cross-encoder, and the top $k$ per language are
returned with full provenance (volume, section, chapter, paragraph index,
and a deep link to the chapter anchor in the hosted edition).

\section{``Ask Vivekananda'': Bridged, Citation-Grounded QA}
\label{sec:qa}

The QA layer turns retrieval into direct answers under two hard
constraints: modern questions must reach era-appropriate material, and
every generated claim must be verifiable. The pipeline has three stages.

\textbf{Stage 1: analyze and bridge.} The question's 15 nearest concept
labels are fetched from the concept index. A single LLM call receives the
question \emph{and this list} and returns (i) 2--4 reformulated
``timeless'' search queries free of modern-only vocabulary, (ii) up to 6
concept labels \emph{copied from the offered list}, and (iii) a one-line
\emph{bridge note} explaining the modern-to-timeless mapping. Restricting
concept choice to labels that verifiably exist in the graph (validated
again against the database after generation) prevents the bridge itself
from hallucinating; if JSON parsing fails, the system degrades to
retrieval with the raw question. The question's language is detected
deterministically from its script, never delegated to the LLM.

\textbf{Stage 2: retrieve.} The original question and each reformulation
are run through the full retrieval stack of
Section~\ref{sec:retrieval}; results are merged, deduplicated per
(language, paragraph), and the pooled candidates are reranked once
against the \emph{original} question. The top 10 passages become the
numbered sources.

\textbf{Stage 3: generate and ground.} A locally hosted
Qwen2.5-7B-Instruct \cite{qwen25} (served via Transformers
\cite{wolf2020transformers}; any OpenAI-compatible endpoint can be
substituted by configuration) receives the numbered sources ---
location-labelled, language-tagged, and trimmed to a window centred on
each passage's best-matching sentence --- and composes a scholarly
third-person answer in the language of the question. The system prompt
forbids claims not supported by the sources and requires a bracketed
citation after every substantive statement. Post-processing then rewrites
each citation marker $[n]$ into a hyperlink into the published online
edition and \emph{deletes} any marker whose index does not correspond to a
retrieved source, so a fabricated citation cannot survive to the reader.
The published editions carry a stable HTML anchor on every paragraph, so
each citation deep-links to the \emph{exact cited paragraph} (falling back
to the chapter anchor where a paragraph anchor is absent) --- the reader
lands on the sentence in question, not merely the top of a chapter. The
interface additionally shows the bridge note,
the reformulated queries, the chosen concepts, and, for every source, the
matched sentence and the full passage --- the entire evidential chain is
one click deep.

\section{System Characteristics}
\label{sec:stats}

\begin{table}[t]
\centering
\caption{Serving-time characteristics (single NVIDIA H200; the system
also runs CPU-only with seconds-level search latency). Model download
sizes excluded.}
\label{tab:perf}
\begin{tabular}{lr}
\toprule
Search, with cross-encoder rerank (mean of 5 queries) & 0.28\,s\\
Search, fusion only & 0.08\,s\\
QA end-to-end (bridge + retrieve + generate) & 6--9\,s\\
Index footprint (7 FAISS indexes + SQLite metadata/graph) & 1.1\,GB\\
Full pipeline build (one GPU) & $<$ 4\,h\\
Incremental rebuild after a source edit & 10--15\,min\\
\bottomrule
\end{tabular}
\end{table}

The stack is deliberately conservative: exact FAISS search, SQLite for
all metadata and graph storage, and a Streamlit interface. There are no
external service dependencies at query time; the system runs entirely
on-premises, which matters both for cost and for institutions that
cannot ship religious-community data to third-party APIs.
Table~\ref{tab:perf} reports measured characteristics. A test suite of
74 CPU-only tests covers the parsers, schema, extraction-output
parsing, fusion, bridging fallbacks, prompt budgeting, citation
linkification, and the concept-graph view.

\section{Evaluation}
\label{sec:eval}

We report three studies: two automatic ones that require no manual
annotation, and a human study of the concept layer. All are reproducible
with evaluation scripts shipped in the
repository.\footnote{\texttt{scripts/eval\_crosslingual.py},
\texttt{scripts/eval\_qa\_integrity.py}, and
\texttt{scripts/eval\_concept\_precision.py}; the raw outputs quoted
here, including the per-annotator judgment files, are included under
\texttt{docs/paper/eval/}.}

\subsection{Known-item cross-lingual retrieval}
\label{sec:eval-knownitem}

The \emph{Vani o Rachana} contains independent Bengali renderings of many
English lectures, which yields chapter-level cross-lingual relevance
judgments at near-zero annotation cost. We identified candidate pairs as
mutual-best matches between the chapter-level embedding indexes (cosine
$\geq 0.70$; 285 candidates), then verified each candidate with the local
LLM shown both chapters' titles and opening paragraphs, retaining
$N = 194$ verified rendered pairs. For each pair, the source chapter's
title and opening paragraph (truncated to 600 characters) is the query,
retrieval is restricted to the \emph{other} language, and the paired
chapter is the sole gold item; a retrieved paragraph counts if it belongs
to the gold chapter. We report Recall@10 and MRR over the top 10 returned
paragraphs in both directions (Table~\ref{tab:knownitem}).

One caveat: pair discovery uses the same embedder as Path A (albeit on
whole-chapter vectors, while queries are titles plus opening paragraphs),
so absolute Path-A numbers may be slightly optimistic; comparisons
\emph{across} configurations are unaffected.

\begin{table}[t]
\centering
\caption{Known-item cross-lingual retrieval over $N=194$ LLM-verified
rendered lecture pairs (EN$\to$BN / BN$\to$EN).}
\label{tab:knownitem}
\begin{tabular}{lcc}
\toprule
Configuration & Recall@10 & MRR\\
\midrule
Path A (direct dense)        & 0.88 / 0.82 & 0.74 / 0.65\\
Path B (concept-mediated)    & 0.14 / 0.12 & 0.07 / 0.04\\
Fused (weighted RRF)         & 0.88 / 0.82 & 0.67 / 0.62\\
Fused + cross-encoder rerank & 0.86 / 0.86 & 0.67 / 0.70\\
\bottomrule
\end{tabular}
\end{table}

Three observations. First, the shared multilingual embedding space alone
is a strong cross-lingual known-item retriever: a query built from an
English lecture finds its Bengali rendering among 14{,}455 Bengali
paragraphs at Recall@10 of 0.88. Second, the concept path is
\emph{not} a known-item retriever and is not meant to be one --- it
expands a query into thematically related material across the whole
graph, which is precisely what known-item metrics penalize; its value
shows in query bridging and in the explanatory ``via concepts'' chains,
not in this table. Third, fusion preserves Path A's recall while adding
the concept candidates, and the cross-encoder markedly improves the
harder BN$\to$EN direction (MRR 0.62$\to$0.70, Recall@10
0.82$\to$0.86) at a small cost in EN$\to$BN, yielding the most balanced
configuration --- the production default.

\subsection{Citation integrity and bridge reliability}
\label{sec:eval-integrity}

We ran the full QA pipeline over a fixed set of 30 questions --- 20
English and 10 Bengali; 15 phrased in deliberately modern vocabulary
(smartphone addiction, social-media envy, job-loss anxiety, burnout) and
15 in timeless vocabulary (fear, concentration, duty, anger) --- and
logged the pipeline's internal decisions (Table~\ref{tab:integrity}).

\begin{table}[t]
\centering
\caption{Citation-integrity and bridging statistics over 30 questions
(20 EN / 10 BN; 15 modern / 15 timeless).}
\label{tab:integrity}
\begin{tabular}{lr}
\toprule
Questions answered (non-empty retrieval) & 30 / 30\\
Stage-1 JSON fallback rate & 0\%\\
Bridge note produced: modern questions & 14 / 15 (93\%)\\
Bridge note produced: timeless questions & 2 / 15 (13\%)\\
Citation markers per answer sentence & 0.47\\
Citation markers emitted (total) & 136\\
Fabricated markers removed by validator & 0 / 136 (0\%)\\
Cross-lingual share of top-10 evidence (bridged) & 32\%\\
Cross-lingual share of top-10 evidence (raw question) & 40\%\\
\bottomrule
\end{tabular}
\end{table}

The bridging stage discriminates almost perfectly between question types:
it produced a bridge note for 93\% of modern questions and only 13\% of
timeless ones, with a 0\% structured-output failure rate. Generated
answers cite densely (one marker per two sentences on average), and in
this run the 7B model emitted \emph{no} citation marker pointing outside
the retrieved source list --- the validator, which deletes such markers,
acted as an idle guardrail rather than an active repair mechanism.

The bridge ablation produced a finding we did not anticipate: the
cross-lingual share of the evidence pool is essentially unchanged by
bridging (32\% bridged vs.\ 40\% raw; 47\% vs.\ 45\% on the modern
subset). Cross-linguality is supplied by the shared embedding space and
the concept graph themselves, not by the reformulation step. The bridge's
measured contributions are therefore selective activation (firing when
and only when the vocabulary gap exists), the era-appropriate
reformulations, and the user-facing explanation --- not an increase in
cross-lingual recall, which the underlying retrieval already provides.

\subsection{Concept-extraction precision: a human study}
\label{sec:eval-precision}

The concept layer is the system's central claim, and it is produced
entirely by an LLM. We therefore ran a human study of its precision.

\textbf{Protocol.} From all 87{,}518 paragraph--concept edges we drew a
uniform random sample of $N=200$ (seed 13). Each item presents one
paragraph, one attached concept label, and the extracted relation.
Annotators judged each item as \emph{correct} (the concept is a correct
reading of the paragraph \emph{at the stated relation}) or
\emph{incorrect} (wrong, merely tangential, or right concept but wrong
relation). Judgments were collected through a small purpose-built web
interface shipped with the repository, which saves after every click and
lets an annotator resume; no annotator saw another's judgments or the
extractor's confidence weight.

\textbf{Annotators.} Three annotators, all familiar with the corpus, each
judged all 200 items independently. Annotators~A and~B are fluent in
both English and Bengali; annotator~C is proficient principally in
Bengali. Since 85 of the 200 sampled items are English, we use C's
judgments only for the Bengali subsample, where that proficiency applies,
and base the whole-sample figures on A and B. All three judgment files
ship with the repository.\footnote{\texttt{docs/paper/eval/annotations/};
the whole-sample figures below are reproduced with
\texttt{eval\_concept\_precision.py compute --users annotator\_a annotator\_b}, and any
other pairing can be recomputed from the same files.}

That scoping is a competence requirement fixed by the task, and the
measured reliability tracks it closely: against~A, annotator~C reaches
$\kappa = 0.49$ on the Bengali items but only $\kappa = 0.14$ --- barely
above chance --- on the English ones. Restricting C to Bengali is
therefore not a way of setting aside an inconvenient judge but of using
each annotator where they are reliable; we report the Bengali
three-annotator panel in full below, and it does not flatter the system.

\textbf{Disclosure.} Annotator~B is an author of this paper; annotators~A
and~C are not. B is marginally the more lenient of A and B (0.705 vs
0.665), so we treat the strict consensus figure --- not B's --- as the
headline number.

\begin{table}[t]
\centering
\caption{Concept-extraction precision over $N=200$ randomly sampled
paragraph--concept edges, annotators A and B. Intervals are Wilson 95\%
(bootstrap for $\kappa$).}
\label{tab:precision}
\begin{tabular}{lcc}
\toprule
Measure & Value & 95\% CI\\
\midrule
Precision, annotator A (non-author)   & 0.665 & [0.597, 0.727]\\
Precision, annotator B (author)       & 0.705 & [0.638, 0.764]\\
\textbf{Precision, consensus} (both correct) & \textbf{0.600} & [0.531, 0.665]\\
Precision, union (either correct)     & 0.770 & [0.707, 0.823]\\
\midrule
Raw agreement                         & 0.830 & ---\\
Cohen's $\kappa$ \cite{cohen1960kappa} & 0.607 & [0.482, 0.720]\\
Gwet's AC$_1$ \cite{gwet2008ac1}       & 0.701 & ---\\
\bottomrule
\end{tabular}
\end{table}

\textbf{Precision is moderate; agreement is substantial.} Consensus
precision is 0.60 --- three in five sampled edges are endorsed by both
judges --- and Cohen's $\kappa$ of 0.61 falls in the ``substantial'' band
on the conventional scale \cite{landis1977kappa}. The residual
disagreement is not systematic: 21 items go one way and 13 the other,
which a McNemar exact test does not distinguish from chance
($p = 0.23$). The two are applying the same rubric with ordinary noise,
which is what an inter-annotator agreement study is supposed to establish
before its precision figure is worth quoting.

\textbf{The extractor's confidence weight is calibrated.} Each
paragraph--concept edge carries a weight the extraction LLM assigns
itself. It predicts human judgment (Spearman $\rho = 0.19$ against
consensus), and thresholding on it works
(Table~\ref{tab:precision-weight}): consensus precision rises from 0.54
below $0.8$ to 0.71 at $0.8$ and above (Fisher exact $p = 0.023$).

\begin{table}[t]
\centering
\caption{Precision by extractor confidence weight. The high-confidence
subgraph is 35.9\% of all edges but still covers 98.1\% of the
paragraphs that carry any concept.}
\label{tab:precision-weight}
\begin{tabular}{lccc}
\toprule
 & Annot.\ A & Annot.\ B & Consensus\\
\midrule
weight $\geq 0.8$ \ ($n=69$)  & 0.783 & 0.812 & \textbf{0.710}\\
weight $< 0.8$ \ ($n=131$)    & 0.603 & 0.649 & 0.542\\
\bottomrule
\end{tabular}
\end{table}

The pruning is nearly free in coverage terms. Applied to the full graph,
weight $\geq 0.8$ retains 31{,}445 of 87{,}518 paragraph--concept edges
(35.9\%) and 3{,}709 of 8{,}362 concepts (44.4\%) --- but
22{,}914 of the 23{,}363 paragraphs that carry any concept
(\textbf{98.1\%}) keep at least one. The low-confidence mass is
redundant attachment to already-covered paragraphs, not the sole
handle on distinct material. A deployment that needs the graph to be
\emph{right} rather than \emph{broad} should therefore threshold at
$0.8$; we report the unpruned graph throughout this paper because Path~B
is used for recall-oriented expansion, where the cost of a wrong concept
is a wasted glance rather than a false statement.

\textbf{Extraction is weaker in Bengali.} Consensus precision over A and
B is 0.68 on English items and 0.54 on Bengali (Fisher exact
$p = 0.043$). The Bengali subsample ($n = 115$) is the one part of the
study carrying a third independent judgment, and the fuller panel
confirms the gap rather than closing it: across A, B and~C the
three-rater Fleiss $\kappa$ is 0.49 and majority-of-three precision is
0.59, still well below the 0.68 measured on English. Adding a
Bengali-proficient annotator thus corroborates the finding with the
language's strongest available judgement.

This is the study's one clearly actionable quality finding beyond the
threshold: the extraction LLM is a predominantly English-trained model
reading Bengali prose, and the concept layer inherits that asymmetry.
Since cross-lingual bridging is the system's central claim, the weaker
half of the graph is precisely the half doing the work of making Bengali
material reachable from English queries --- a Bengali-strong extractor,
or a second extraction pass over Bengali paragraphs alone, is the most
valuable single upgrade available to the pipeline.

\textbf{Relations.} The dominant \emph{discusses} (87.5\% of the sample)
reaches consensus precision 0.62, while the two rare relations are worse
(\emph{exemplifies} 0.47, $n=17$; \emph{contrasts} 0.50, $n=8$) --- on
samples too small to support a strong claim, but consistent with the
intuition that a directional judgment is harder for the extractor than a
topical one.

\textbf{On the rubric.} The binary judgment conflates two readings of
``the concept describes the paragraph'': whether the paragraph \emph{is
about} the concept, and whether the concept is merely \emph{present} in
it. For a S\=utra gloss such as ``The success of Yogis differs according
as the means they adopt are mild, medium, or intense'' tagged
\texttt{success}, or a Bengali travel note announcing a forthcoming
Vedanta lecture tagged \texttt{vedanta}, the first reading rejects and
the second accepts. Both annotators converged on the more permissive
presence reading, so the 0.60 we report should be understood as a
precision under that reading; a study that fixed the stricter aboutness
criterion for every annotator would return a lower number. A revised
study should use a three-point scale (\emph{central} / \emph{mentioned} /
\emph{unsupported}) rather than force this distinction into a binary, and
should fix the reading in the guidelines instead of leaving it to be
resolved per annotator.

\textbf{What the errors look like.} The characteristic failure is not a
bizarre concept but a plausible one attached to a paragraph that does not
support it. Narrative, epistolary, and connective passages attract the
corpus's thematic vocabulary regardless of content: a Bengali letter
remarking that a book ``has been well received here, but the publishers
do not seem to be pushing sales'' (our translation) is tagged
\texttt{knowledge-realization}; a subscription appeal for a memorial is
tagged \texttt{liberation}. The
extractor appears to condition on the corpus's overall register as much
as on the paragraph in front of it. This is precisely the failure mode
that the confidence threshold suppresses.

\textbf{Implication for the system's claims.} The concept layer should
be understood as a \emph{recall-oriented, imperfect expansion and
explanation layer}, not a curated ontology. This is consistent with
Table~\ref{tab:knownitem}, where Path~B is a poor known-item retriever
on its own; its contribution is cross-lingual reach and the
``via concepts'' explanations, both of which degrade gracefully under
error because every retrieved passage is displayed to the reader with
its source. A wrong concept costs an irrelevant result, not a false
assertion about Vivekananda --- the citation discipline of
Section~\ref{sec:qa}, not the concept graph, is what bounds that risk.


\section{Case Studies}
\label{sec:cases}

We illustrate the bridging behaviour with unedited system outputs.

\textbf{Modern question, English.} For \emph{``How do I get rid of
mobile phone addiction?''} the analysis stage produced the bridge note
\emph{``Mobile-phone addiction is a modern form of attachment of the
senses and a habit that enslaves the mind''}, selected the corpus
concepts \texttt{excessive-attachment} and \texttt{self-control}, and
reformulated the question into queries such as \emph{``How can I
overcome excessive attachment to sense objects?''} and \emph{``What is
the remedy for habit that enslaves the mind?''}. Retrieval returned ten
passages spanning \emph{Inspired Talks}, the Raja-Yoga lectures, and ---
notably --- Bengali chapters on liberation (\emph{mukti}) and
sense-withdrawal (\emph{pratyahara}), which no lexical and no unbridged
dense query for ``mobile phone'' reaches. The generated answer opens by
restating the bridge, then prescribes non-attachment and graded practice
of sense-withdrawal, with every paragraph citing its sources.

\textbf{Modern question, second example.} For \emph{``How should one
deal with social media envy?''} the bridge note read \emph{``Social
media envy is a modern form of coveting the social status and approval
of others\ldots''} --- a mapping to covetousness and approval-seeking
that the authors did not anticipate but that the concept inventory
supported.

\textbf{Bengali question.} A Bengali question on freedom from mobile
addiction (\emph{``mobile asakti theke muktir upay ki?''}) was detected
as Bengali, bridged through the same English-labelled concepts, answered
\emph{in Bengali}, and grounded predominantly in Bengali sources while
still drawing on English lectures --- the concept layer making the
cross-lingual evidence pool available in both directions.

\section{Limitations and Ethical Considerations}
\label{sec:limitations}

\textbf{An imperfect concept layer.} The human study of
Section~\ref{sec:eval-precision} reports strict-consensus precision of
0.60 over paragraph--concept edges, and 0.54 on Bengali material. The
concept graph is a useful expansion and explanation layer, but it is not
a scholarly index and should not be cited as one. Users who need higher
precision should threshold at extractor weight $\geq 0.8$, which lifts
consensus precision to 0.71 at almost no cost in paragraph coverage.

\textbf{Limits of that study.} The whole-sample figure rests on two
annotators, one of them an author of this paper and marginally the more
lenient of the two, judging 200 items --- enough to place the precision
within roughly $\pm 0.07$, but not to support fine comparisons. A third
annotator, proficient principally in Bengali, contributes only to the
Bengali subsample; the English half therefore carries no judgement from
outside the author's immediate circle. The binary rubric conflated ``the
paragraph is about this concept'' with ``this concept appears in the
paragraph'' and left the choice to each annotator; the two whole-sample
annotators settled on the more permissive reading, so the headline figure
is the optimistic end of the plausible range --- the Bengali-proficient
annotator, who judged more strictly, would place it lower. A three-point
scale, guidelines that fix the reading, and a larger and wholly
independent annotator pool are all needed before the precision figure can
be treated as settled rather than indicative.

\textbf{Answer faithfulness is still unevaluated.} Citation integrity is
measured structurally (marker validity), not semantically: we verify that
every citation points at a retrieved passage, not that the passage
supports the claim attached to it. Known-item judgments likewise come
from rendered lecture pairs rather than topical relevance assessment. A
scholar-annotated faithfulness study of generated answers remains the
most important piece of future work.

\textbf{Generation quality floor.} The default 7B answer model
occasionally paraphrases loosely when quoting, and may present material
from a Bengali source in English without consistently marking it as
translated, despite prompt instructions. The citation mechanics bound
the damage --- the linked passage is always one click away --- but do not
eliminate it. Configurable substitution of a stronger model raises the
floor.

\textbf{Extraction bias.} The concept inventory reflects the judgement
of the extraction LLM; systematic blind spots in that model's reading of
Vedantic material would propagate to the graph. The alias tables and
per-paragraph provenance make audits feasible.

\textbf{Cultural sensitivity.} The corpus is scripture-adjacent for a
living community. We consider verifiable grounding, the visible bridge
explanation, and on-premises deployment to be ethical requirements for
this material, not merely engineering choices, and we would caution
against deploying fluent-generation systems over such corpora without
citation enforcement of comparable strictness.

\section{Conclusion}

\system{} demonstrates that a modest, fully local pipeline --- one
extraction pass by a mid-sized LLM, a shared multilingual embedding
space, and a canonical-label concept graph --- suffices to make a
bilingual, non-parallel classical corpus searchable by meaning,
answerable in either language, and honest about its evidence. The
temporal bridge turns the system's most obvious weakness (a
nineteenth-century vocabulary) into an explicit, inspectable mapping
rather than a silent retrieval failure. The architecture assumes nothing
Vivekananda-specific: any corpus with stable canonical structure and one
or more languages covered by the embedding model --- the Ramakrishna
\emph{Kathamrita} literature, Gandhi's collected works in English,
Gujarati and Hindi, or patristic corpora in Greek and translation ---
could be indexed with the same pipeline. The human study also marks the
clearest directions for the next iteration. At strict-consensus precision
of 0.60 the concept layer is broad but imperfect, and two of its
properties point at cheap gains: the extractor's own confidence weight
separates good edges from bad ones well enough that a calibrated
threshold buys accuracy almost for free, and the precision gap between
English (0.68) and Bengali (0.54) says the largest single quality deficit
is not the pipeline but the extraction model's weaker reading of Bengali
--- the very half of the corpus that cross-lingual bridging exists to
open up. A Bengali-capable extractor and community-driven correction
through the audit trail the system already exposes are therefore the
first steps; graph-embedding refinement of the concept layer and a
scholar-annotated faithfulness study of generated answers are the
further ones.

\section*{Acknowledgements}

Facilities utilized in this research were supported by the Fund for
Improvement of S\&T Infrastructure (FIST) program of the Department of
Science and Technology (DST), India [Sanction No.: SR/FST/MS-I/2022/116].

The author thanks the volunteer annotators who each independently judged
all 200 items of the concept-precision study reported in
Section~\ref{sec:eval-precision}.

\begin{thebibliography}{99}\setlength{\itemsep}{2pt}

\bibitem{lewis2020rag}
P.~Lewis, E.~Perez, A.~Piktus, et al.
\newblock Retrieval-augmented generation for knowledge-intensive NLP tasks.
\newblock In \emph{Advances in Neural Information Processing Systems 33
(NeurIPS)}, 2020.

\bibitem{karpukhin2020dpr}
V.~Karpukhin, B.~O\u{g}uz, S.~Min, et al.
\newblock Dense passage retrieval for open-domain question answering.
\newblock In \emph{Proceedings of EMNLP}, 2020.

\bibitem{gao2023ragsurvey}
Y.~Gao, Y.~Xiong, X.~Gao, et al.
\newblock Retrieval-augmented generation for large language models: A survey.
\newblock \emph{arXiv preprint arXiv:2312.10997}, 2023.

\bibitem{edge2024graphrag}
D.~Edge, H.~Trinh, N.~Cheng, et al.
\newblock From local to global: A graph RAG approach to query-focused
summarization.
\newblock \emph{arXiv preprint arXiv:2404.16130}, 2024.

\bibitem{chen2024bgem3}
J.~Chen, S.~Xiao, P.~Zhang, K.~Luo, D.~Lian, and Z.~Liu.
\newblock BGE M3-Embedding: Multi-lingual, multi-functionality,
multi-granularity text embeddings through self-knowledge distillation.
\newblock \emph{arXiv preprint arXiv:2402.03216}, 2024.

\bibitem{reimers2019sbert}
N.~Reimers and I.~Gurevych.
\newblock Sentence-BERT: Sentence embeddings using Siamese BERT-networks.
\newblock In \emph{Proceedings of EMNLP-IJCNLP}, 2019.

\bibitem{johnson2019faiss}
J.~Johnson, M.~Douze, and H.~J\'egou.
\newblock Billion-scale similarity search with GPUs.
\newblock \emph{IEEE Transactions on Big Data}, 7(3):535--547, 2019.

\bibitem{cormack2009rrf}
G.~V. Cormack, C.~L.~A. Clarke, and S.~B\"uttcher.
\newblock Reciprocal rank fusion outperforms Condorcet and individual rank
learning methods.
\newblock In \emph{Proceedings of SIGIR}, 2009.

\bibitem{qwen25}
Qwen Team.
\newblock Qwen2.5 technical report.
\newblock \emph{arXiv preprint arXiv:2412.15115}, 2024.

\bibitem{kwon2023vllm}
W.~Kwon, Z.~Li, S.~Zhuang, et al.
\newblock Efficient memory management for large language model serving with
PagedAttention.
\newblock In \emph{Proceedings of the 29th ACM Symposium on Operating Systems
Principles (SOSP)}, 2023.

\bibitem{wolf2020transformers}
T.~Wolf, L.~Debut, V.~Sanh, et al.
\newblock Transformers: State-of-the-art natural language processing.
\newblock In \emph{Proceedings of EMNLP: System Demonstrations}, 2020.

\bibitem{bird2004nltk}
S.~Bird and E.~Loper.
\newblock NLTK: The natural language toolkit.
\newblock In \emph{Proceedings of the ACL Interactive Poster and
Demonstration Sessions}, 2004.

\bibitem{ji2023survey}
Z.~Ji, N.~Lee, R.~Frieske, et al.
\newblock Survey of hallucination in natural language generation.
\newblock \emph{ACM Computing Surveys}, 55(12):1--38, 2023.

\bibitem{cohen1960kappa}
J.~Cohen.
\newblock A coefficient of agreement for nominal scales.
\newblock \emph{Educational and Psychological Measurement}, 20(1):37--46, 1960.

\bibitem{landis1977kappa}
J.~R. Landis and G.~G. Koch.
\newblock The measurement of observer agreement for categorical data.
\newblock \emph{Biometrics}, 33(1):159--174, 1977.

\bibitem{gwet2008ac1}
K.~L. Gwet.
\newblock Computing inter-rater reliability and its variance in the presence
of high agreement.
\newblock \emph{British Journal of Mathematical and Statistical Psychology},
61(1):29--48, 2008.

\bibitem{malhas2020ayatec}
R.~Malhas and T.~Elsayed.
\newblock AyaTEC: Building a reusable verse-based test collection for Arabic
question answering on the Holy Qur'an.
\newblock \emph{ACM Transactions on Asian and Low-Resource Language
Information Processing}, 19(6):1--21, 2020.

\end{thebibliography}

\end{document}
